\chapter{Future Work: Predicting Numerical Variability}
\label{chap:future-work}

\section{Motivation and main question}

The first phase of this research demonstrated that numerical variability (NV) in functional MRI graph measures is measurable and can influence statistical conclusions. These findings establish the importance of evaluating numerical stability, particularly when the observed effects are small or close to a decision threshold. However, direct assessment with Monte Carlo Arithmetic (MCA) requires repeated executions of the processing workflow and is therefore computationally expensive. In addition, most researchers do not routinely have access to an MCA-instrumented environment. Consequently, numerical stability is rarely evaluated as part of a conventional neuroimaging analysis.

The next phase of this research is motivated by a practical question:

\begin{quote}
\emph{Can the numerical variability of graph measures be predicted from the functional connectome produced by a single standard execution?}
\end{quote}

Here, a standard or \emph{IEEE} execution refers to ordinary processing without MCA perturbations. If the resulting functional connectivity (FC) matrix contains information about its own numerical sensitivity, it may be possible to estimate NV without repeating the complete workflow. Such an estimate would not replace MCA as the reference method for characterizing numerical uncertainty, but it could provide an accessible screening tool for identifying potentially unstable results.

\section{Why prediction may be possible}

Preliminary exploratory analyses investigated how MCA perturbations in FC propagate into graph-metric NV. The findings suggest that this propagation is structured rather than uniform. In particular, connections whose correlation values are close to the graph-construction threshold are more likely to switch between being included and excluded. These edge changes can alter network topology and amplify small numerical perturbations into larger differences in graph measures.

Threshold proximity alone does not fully determine this effect. The spatial and topological locations of the affected edges are also important: a perturbation involving a highly connected or otherwise influential part of the network may have a different consequence from an equally sized perturbation elsewhere. The analysis further indicated that binary graph variability arises through changes in topology, whereas weighted graph variability can reflect both changing topology and changing connection weights. Together, these findings suggest that the standard connectome contains potentially informative features---including connection magnitude, distance from the threshold, and the organization of vulnerable edges---that may indicate future numerical sensitivity. This provides the main rationale for prediction.

\section{Preliminary prediction results}

Initial experiments used a simple BrainNetCNN-style model to predict graph-measure NV from the IEEE connectivity matrix. This model was selected because its edge-to-edge and edge-to-node operations are designed to retain relationships among connections and brain regions rather than treating connectivity values as unrelated features. A logarithmic transformation of NV was used to reduce the influence of its strongly skewed distribution.

The results provide preliminary evidence that prediction may be feasible, although performance is not consistent across all graph measures or thresholds. Some outcomes, including degree and eigenvector centrality in lower-threshold graphs and betweenness or clustering in other threshold settings, showed more promising prediction than others. 

For several graph measures, the Spearman correlations between predicted and observed NV were informative even when prediction of the exact NV magnitude remained limited. This indicates that the model may already contain useful information for ordering participants from less to more numerically sensitive. Such ordering is relevant in practice because it could help identify participants for whom a direct MCA analysis should be prioritized.

These findings remain preliminary and do not yet establish a generally accurate predictor. Their importance is instead that they support the central premise of this phase: the IEEE connectome appears to contain information about numerical fragility, but the usefulness of that information differs across graph measures and graph-construction choices.

\section{Challenges and proposed direction}

Predicting NV presents several conceptual challenges. The response to MCA perturbations is itself random, and NV estimated from a finite number of realizations is therefore uncertain. This uncertainty can limit the apparent accuracy of a model even when a relationship with the underlying numerical sensitivity exists. In addition, the IEEE connectome may not contain all relevant information about how a processing workflow responds to perturbation. Numerical sensitivity may also depend on intermediate computations or computational conditions that are not represented in the final FC matrix.

The next stage will investigate whether prediction can be improved. However, exact numerical prediction is not the only potentially useful outcome. If the uncertainty of the target or missing sensitivity information places a practical limit on magnitude prediction, the work will evaluate whether participants can instead be ranked reliably by NV or classified as having comparatively high numerical sensitivity. Identifying highly sensitive participants would still provide a practical basis for directing computationally expensive MCA analyses toward the cases in which they are most needed.

\section{Expected contribution}

This phase will determine whether information available from a single standard connectome can be used to anticipate the numerical fragility of derived graph measures. It will connect the mechanistic evidence on threshold crossing and network topology with a practical prediction problem. A successful approach could make numerical-stability assessment more accessible by providing either an estimate of NV or a reliable indication of which participants are most likely to be sensitive. If prediction remains limited, the results will still clarify which aspects of numerical sensitivity are absent from the standard connectome and where direct perturbation-based evaluation remains necessary.
